\documentclass[11pt]{amsart}

\usepackage[T1]{fontenc}
\usepackage{lmodern}
\usepackage{microtype}
\usepackage{amsmath,amssymb}
\usepackage{booktabs}
\usepackage[hidelinks]{hyperref}

\hypersetup{
  pdftitle={Conjecture 1.7 of Qu and Zang Is False}
}

\newtheorem{theorem}{Theorem}
\theoremstyle{remark}
\newtheorem{remark}[theorem]{Remark}

\newcommand{\Podd}{\mathcal P_{\mathrm{odd}}}
\newcommand{\ind}{\mathbf 1}

\title[Conjecture 1.7 of Qu and Zang]
{Conjecture 1.7 of Qu and Zang Is False}

\date{}

\subjclass[2020]{05A17, 05A20, 11P81}
\keywords{integer partitions, hook lengths, regular partitions,
partition inequalities}

\begin{document}

\begin{abstract}
Let $b_{2,j}(n)$ be the total number of hooks of length $j$ among the
$2$-regular partitions of $n$.  We prove that, for every odd $n\ge 7$,
\[
 b_{2,n-3}(n)=5,
 \qquad
 b_{2,n-2}(n)=\frac{n-1}{2}.
\]
Hence, for every even $k\ge 10$,
\[
 b_{2,k}(k+3)=5<\frac{k+2}{2}=b_{2,k+1}(k+3),
\]
which disproves Conjecture~1.7 of Qu and Zang.
\end{abstract}

\maketitle

A $t$-regular partition is a partition with no part divisible by $t$.
For $t\ge2$ and $j\ge1$, let $b_{t,j}(n)$ denote the total number of
hooks of length $j$ among all $t$-regular partitions of $n$.  Singh and
Barman conjectured that
\[
 b_{2,k}(n)\ge b_{2,k+1}(n)
 \qquad (k\ge3,\ n\ge0,\ n\ne k+1)
\]
\cite[Conjecture~1.6]{SinghBarman}.  Qu and Zang proved that this fails
for every odd $k$, proved the cases $k=4$ and $k=6$, and conjectured the
inequality for every even $k\ge8$ \cite[Theorems~1.4--1.6 and
Conjecture~1.7]{QuZang}.  The following elementary formulas disprove
that conjecture for every even $k\ge10$.

\begin{theorem}\label{thm:main}
For every odd integer $n\ge7$,
\[
 b_{2,n-3}(n)=5,
 \qquad
 b_{2,n-2}(n)=\frac{n-1}{2}.
\]
Consequently, for every even integer $k\ge4$,
\[
 b_{2,k}(k+3)-b_{2,k+1}(k+3)=\frac{8-k}{2}.
\]
In particular, the conjectured inequality fails for every even
$k\ge10$, and Conjecture~1.7 of Qu and Zang is false.
\end{theorem}

\begin{proof}
A $2$-regular partition is a partition into odd parts.  Let
$\lambda\in\Podd(n)$, written as
$\lambda=(\lambda_1,\ldots,\lambda_\ell)$.  For a cell $(r,c)$ of its
Young diagram, define the defect by
$d_\lambda(r,c)=n-h_\lambda(r,c)$.  Counting the cells outside the hook
of $(r,c)$ gives
\begin{equation}\label{eq:defect}
 d_\lambda(r,c)
 =\sum_{i<r}\lambda_i+(c-1)
  +\sum_{i>r}\bigl(\lambda_i-\ind_{\{\lambda_i\ge c\}}\bigr).
\end{equation}
We claim that the complete lists of cells of defects $2$ and $3$ are as
follows; repeated parts are written in exponential notation.
\begin{center}
\begin{tabular}{ccl}
\toprule
Defect & Partition $\lambda$ & Cell $(r,c)$ \\
\midrule
$2$ & $(n)$ & $(1,3)$ \\
    & $(1^n)$ & $(3,1)$ \\
    & $(j,3,1^{\,n-j-3})$, $j\in\{3,5,\ldots,n-4\}$ & $(1,1)$ \\
\midrule
$3$ & $(n)$ & $(1,4)$ \\
    & $(n-2,1^2)$ & $(1,2)$ \\
    & $(3,3,1^{\,n-6})$ & $(2,1)$ \\
    & $(3,1^{\,n-3})$ & $(2,1)$ \\
    & $(1^n)$ & $(4,1)$ \\
\bottomrule
\end{tabular}
\end{center}

We verify exhaustiveness.  First suppose $r=1$.  Equation
\eqref{eq:defect} implies $c\le d_\lambda(1,c)+1$.  For defect $2$, the
cases $c=3,2,1$ give, respectively, no lower row; exactly one lower part
equal to $1$, which would make $\lambda_1=n-1$ even; or
\[
 \sum_{i>1}(\lambda_i-1)=2.
\]
The last equation forces exactly one lower part equal to $3$, with all
remaining lower parts equal to $1$.  This gives $(j,3,1^{n-j-3})$.
Here $j$ is odd and $n-j-3$ is a nonnegative odd integer, so
$j\in\{3,5,\ldots,n-4\}$.

For defect $3$, the cases $c=4,3,2,1$ give, respectively, no lower row;
exactly one lower part equal to $1$, again making $\lambda_1$ even;
either two lower parts equal to $1$ or one lower part equal to $3$, of
which only the former makes $\lambda_1$ odd; or an impossibility, since
$\sum_{i>1}(\lambda_i-1)$ is even.  These are exactly the listed cells
with $r=1$.

Now suppose $r\ge2$.  If $c\ge2$, then $\lambda_1\ge\lambda_r\ge3$, so
\[
 \sum_{i<r}\lambda_i+(c-1)\ge\lambda_1+1\ge4.
\]
Thus a cell of defect at most $3$ must have $c=1$.  Set
\[
 U=\sum_{i<r}\lambda_i,
 \qquad
 L=\sum_{i>r}(\lambda_i-1).
\]
Then $d_\lambda(r,1)=U+L$, where $L$ is even and
\[
 U\ge r-1,
 \qquad
 U\equiv r-1\pmod 2.
\]
If the defect is $2$, these conditions force $r=3$, $U=2$, and $L=0$;
hence $\lambda=(1^n)$ and $(r,c)=(3,1)$.

If the defect is $3$, then $U$ is odd, so $r$ is even, and
$r-1\le U\le3$ gives $r\in\{2,4\}$.  For $r=2$, if
$\lambda_1=1$, then monotonicity gives $L=0$ and the defect is $1$.
Thus $\lambda_1=3$, which forces $L=0$ and
$\lambda_2\in\{1,3\}$, giving
$(3,1^{n-3})$ and $(3,3,1^{n-6})$.  For $r=4$, one has $U=3$ and
$L=0$, so $\lambda=(1^n)$.  This proves the classification.

There are five partition--cell pairs of defect $3$.  For defect $2$,
there are two isolated pairs and $(n-5)/2$ admissible values of $j$.
Therefore
\[
 b_{2,n-3}(n)=5,
 \qquad
 b_{2,n-2}(n)=2+\frac{n-5}{2}=\frac{n-1}{2}.
\]
Substituting $n=k+3$ proves the remaining assertions.
\end{proof}

\begin{remark}
Taking $k=10$ gives the explicit counterexample
\[
 b_{2,10}(13)=5<6=b_{2,11}(13).
\]
In the arXiv version of \cite{QuZang} (arXiv:2501.13753v1, p.~3),
Conjecture~1.7 is introduced by the sentence ``After checking
$8\le k\le20$ and $n\le10000$, we propose the following
conjecture\dots''; the pair $(k,n)=(10,13)$ lies in the range so
described.  At $k=8$,
Theorem~\ref{thm:main} gives equality at $n=11$.  Together with the
results of Qu and Zang, $k=8$ is the only fixed parameter $k\ge3$ not
decided by these results.
\end{remark}

\begin{thebibliography}{9}

\bibitem{QuZang}
W. Qu and W. J. T. Zang,
\emph{On the hook length biases of the $2$- and $3$-regular partitions},
J. Number Theory \textbf{280} (2026), 455--480.
\href{https://doi.org/10.1016/j.jnt.2025.08.016}
{doi:10.1016/j.jnt.2025.08.016};
\href{https://arxiv.org/abs/2501.13753}{arXiv:2501.13753}.

\bibitem{SinghBarman}
G. Singh and R. Barman,
\emph{Hook length biases in ordinary and $t$-regular partitions},
J. Number Theory \textbf{264} (2024), 41--58.
\href{https://doi.org/10.1016/j.jnt.2024.05.001}
{doi:10.1016/j.jnt.2024.05.001}.

\end{thebibliography}

\end{document}